\documentclass[11pt,a4paper]{article}
\usepackage[margin=21mm]{geometry}
\usepackage{amsmath,amssymb,amsthm}
\setlength{\emergencystretch}{2em}
\newtheorem{lemma}{Lemma}
\newtheorem{theorem}[lemma]{Theorem}
\newcommand{\gen}[1]{\langle#1\rangle}
\title{Pair blocks are incompatible with the critical head}
\author{Human Compression project: L044 continuation}
\date{}
\begin{document}\maketitle
Retain the complete critical-head hypotheses: $n=2r\ge18$,
$G\le\operatorname{Alt}(\Omega)$ is transitive with connected
string C-marking, and every full ambient maximal omission is
intransitive. The acquired six-point head has support $D$.
Write
\[
 E=\Omega\setminus D,\quad B=G_{\ge4},\quad K=G_{\ge5},
 \quad\sigma=\rho_4,\quad x=b^\sigma\ne b.
\]
The acquired early-cut calculation gives an odd involution
$\delta$ supported on $E$, fixing $b,x$ and centralizing $K$.
The group $B$ fixes $D$ pointwise, and $K$ fixes $b$.

\begin{lemma}[Two alternatives at rank two below the degree]
\label{row}
Let $H=\gen{h_1,\ldots,h_{N-2}}$ be a faithful string
C-group on $N\ge5$ points, with connected noncommutation
diagram. Suppose deleting each mark changes its orbit partition.
If $H$ is transitive, its centralizer in $S_N$ is trivial.
If it is intransitive, it fixes one point and is natural
$S_{N-1}$ on the rest; its marks are the single edge
transpositions of a path, in path order or reverse order.
\end{lemma}
\begin{proof}
In the transitive case apply the complete accepted L047 theorem.
For odd $N$, rank $N-2>N/2$ and all its proper exceptions
have even degree. For even $N\ge8$,
$N-2>N/2+1$, greater than every proper exception rank.
Thus in these degrees $H$ is natural $A_N$ or $S_N$,
whose centralizer is trivial. When $N=6$, the proper
rank-four imprimitive exceptions have types, up to reversal,
\[
 (3,2,3),\ (2,3,3),\ (2,3,4),\ (2,4,3).
\]
All have disconnected noncommutation diagrams. The only
connected proper possibility is the primitive degree-six
$S_5$, also with trivial centralizer. Indeed, in a faithful
primitive nonregular action a point stabilizer is maximal
and self-normalizing: its normalizer is either itself or
the whole group, and the latter would make it normal and
then trivial by faithfulness. A centralizing permutation
is determined by the image of a point, and the possible
images correspond to the normalizer quotient. This proves
the asserted triviality in every transitive case.

If $H$ is intransitive, choose for each mark an edge joining
different orbits of its omission. The $N-2$ selected edges
form a forest, since a cycle would supply an avoiding path.
There are exactly two forest components; they are precisely
the two $H$-orbits, since $H$ is intransitive.
Deleting a selected edge $j$ splits its tree into two pieces.
The actual omission contains every other selected edge and
cannot reconnect these pieces, so every other mark preserves
that cut. Within a tree the only two distinct vertices not
separated by a cut other than cut $i$ are the endpoints of
edge $i$. Thus each $h_i$ is exactly its selected
transposition. Each tree supplies its full natural symmetric
group. If both trees had edges, their two nonempty sets of
marks would commute, disconnecting the diagram. One tree
is therefore a singleton. Finally, a vertex of degree three
in the other tree would give a triangle of noncommutation
among its incident edge transpositions, impossible in a string.
That tree is a path, and its connected string ordering is
its path order or its reverse.
\end{proof}

\begin{theorem}
Under the standing hypotheses, $B$ cannot be transitive on $E$
and preserve a partition into pairs. This covers every
critical degree $n\ge18$ and uses no alternating rank induction.
\end{theorem}
\begin{proof}
Suppose there are $m=(n-6)/2=r-3\ge6$ pairs. The rank
of $B$ is $d=r-4=m-1$. Each distinguished omission in $B$
is intransitive on $E$: otherwise adjoining the head and
$\rho_3$, which links the head to $b$, would make the
corresponding full ambient omission transitive.

Let $M$ be a minimal distinguished subset generating the
whole action on the pairs, and let $c$ count the actual
point orbits of $\gen M$ on $E$. Transitivity on the pairs
makes every point orbit meet every pair, so $c\le2$.
The acquired actual-orbit forest and Whiston bound give
\[
               d-|M|\le c-1\le1,\qquad |M|\le m-1.
\]
Minimality makes the block images of $M$ independent.
All marks other than $\sigma$ fix the pair containing $b$.
Hence $\sigma\in M$ and moves that pair. Let $b'$ be its
other point; then $x$ is outside $\{b,b'\}$.

If $|M|=m-1$, all distinguished block images are independent.
The $m-2$ images from $K$ fix the pair $\{b,b'\}$ and
attain maximum independent length on the other $m-1$ pairs.
Whiston's equality clause makes their image natural $S_{m-1}$.
The group $K$ fixes $b,b'$ pointwise. On the other points
it has either one transitive orbit or two orbits, each meeting
each pair once. In the first case, centralization and fixation
of $x$ make $\delta$ fix the whole orbit. In the second,
the two row actions are faithful natural $S_{m-1}$ actions;
$\delta$ fixes the row containing $x$ pointwise, and the
other row is fixed by its trivial natural centralizer.
The remaining fixed set is $\{b,b'\}$; fixation of $b$
forces fixation of $b'$. Thus $\delta=1$, a contradiction.

We may therefore suppose $|M|=m-2$. The forest inequality
forces $c=2$. Let $U,V$ be the two point orbits, with
$b\in U$ and $b'\in V$. Each meets every pair once,
so either row action of $\gen M$ is faithful: an element
fixing one row fixes every pair and preserves both rows,
hence fixes the other row as well. It inherits full IP
and all $m-2$ nonidentity distinguished marks.

Its noncommutation diagram is connected. A component not
containing $\sigma$ fixes $b$ and commutes with the other
components. Transitivity then makes it fix the entire row,
contrary to faithfulness. Since $M$ contains the first mark
of $B$ and omits just one mark, connectedness forces it
to omit the last mark: an internal missing index would
separate two nonempty commuting intervals. Thus
\[
 M=\{\rho_4,\ldots,\rho_{r-2}\},\qquad
 H=\gen{M\setminus\{\sigma\}}=G_{5,\ldots,r-2}.
\]
The group $H$ fixes $b,b'$ and acts faithfully and diagonally
on $U\setminus\{b\}$ and $V\setminus\{b'\}$,
each of degree $N=m-1\ge5$, with rank $m-3=N-2$.
Its diagram is connected, being an interval of the connected
marking of $M$.

Every row omission in $H$ changes its orbit partition.
Otherwise the identical row partitions on both sides, with
$b,b',D$ fixed, would leave the actual point-orbit partition
of $H$ unchanged. Adjoining all the other ambient marks
would give the same generated orbit equivalence relation
with or without the deleted mark. The full ambient omission
would be transitive, a contradiction. Thus Lemma~\ref{row}
applies to the faithful row marking of $H$.

If the row action of $H$ is transitive, $\delta$ centralizes
it and fixes $x\in U\setminus\{b\}$; here $x$ lies
in $U$ because $\sigma\in\gen M$ preserves $U$.
It fixes that whole row, and its action on the other row
is trivial by Lemma~\ref{row}. The fixed points $b,b'$
are fixed as above, again giving $\delta=1$.

If the row action is intransitive, write $p$ for its unique
singleton in $U\setminus\{b\}$ and $W=U\setminus\{b,p\}$.
On $W$, of size $m-2\ge4$, its marks are the transpositions
of a path. Let $e=\rho_5$ be its first edge and let $a$
be the endpoint detached by deleting $e$. Then
\[
        \operatorname{Fix}_U(G_{6,\ldots,r-2})=\{b,p,a\}.
\]
The mark $\sigma$ commutes with this entire subgroup and
therefore preserves that three-point fixed set. On the other
hand $\gen M=\gen{\sigma,H}$ is transitive on $U$,
while $H$ fixes both $b,p$. Since $\sigma$ is an involution,
transitivity forces it to send both $b$ and $p$ into $W$:
a fixed point or an interchange $b\leftrightarrow p$ would
leave a singleton or the pair $\{b,p\}$ invariant.
Both images would consequently have to equal $a$, impossible
for a permutation. This excludes the last case.
\end{proof}

\paragraph{Attribution and boundary.}
The critical-head calculation and actual-orbit forest/Whiston
interfaces remain acquired. The full accepted L047 theorem,
including the exact small exception list and its common theory,
finite certificate, software and distinct reviews, remains paid.
The new step applies it to the endpoint omission on each literal
row; a saturated forest then excludes the intransitive alternative.
No B5 classification, smaller-degree alternating induction,
pending packet or new computation is used. Other block systems,
opposite heads, the unshared endpoint and the complete L044
recovery remain separate; no whole-original or question credit
is claimed.
\end{document}
